\section{异构运输调度与候选方案权衡}
\subsection{时限与资源模型}
Q2 在同一物理模型上增加实体无人机、电池及配送时限。定义 $x_{pb}=1$ 表示货箱 $b$ 由架次 $p$ 承运，且 $\sum_px_{pb}=1$。机型、访问顺序确定后，公共航段引擎给出作业时长 $t_p$、能耗 $E_p$ 和货箱交付偏移 $\delta_{pb}$，因而
\begin{equation}
C_b=s_p+\delta_{pb},\quad
\ell_b\ge C_b-e_b,\quad\ell_b\ge0,\quad
C_b\le d_b\quad(b\in\mathcal B_{\rm hard}).
\label{eq:deadline}
\end{equation}
$e_b$ 是期望时间，$d_b$ 是硬截止时间：医疗物资以期望时间作为硬约束，首批保障以首批截止为硬约束，同时满足两种身份时取更早者。其他物资期望时间进入加权延误 $W=\sum_bw_b\ell_b$。

若架次 $p,p'$ 共用无人机，则区间 $[s_p,s_p+t_p)$ 不得重叠；若共用电池，则区间延伸至 $s_p+t_p+t_{\rm chg}(1-E_p/E_g)$。把电池与无人机绑定成不可拆的固定组合会丢失共享电池的调度能力，因此分别记录两类分配变量。多指标为
\begin{equation}
\min\left(W,\ T_2,\ E_2,\ N_2\right),\qquad
T_2=\max_p(s_p+t_p),\quad E_2=\sum_pE_p.
\label{eq:q2objective}
\end{equation}
式\eqref{eq:q2objective}表达权衡关系；实现中通过有限候选搜索和人工明确选点处理，并未求出完整 Pareto 前沿。

\subsection{确定性启发式与选点原则}
首先以紧迫货箱为锚构造满足硬时限的路线，优先保障硬时限，再综合延误、承运质量和箱数。其次尝试合并架次，枚举候选访问顺序并重新选择可用资源、重算出发时间；每条路线最多枚举五个服务区，以限制阶乘复杂度。每次合并后的载荷、逐箱完成时间、能源和资源占用都重新计算。

保存架次数从 24 至 17 的八个候选，而不是只留下选中的结果。附录表\ref{tab:frontier}与图\ref{fig:optimization}表明，22 架次相较更少架次保留了较短完工历时和较低延误；相比 23 架次，它降低能耗，但增加部分软时限延误。因此选用 \QTwoTrips{} 架次是一个显式折中，不能称为由所有目标共同唯一决定的最优点。
\paperfig{fig05_optimization}{较少架次需要付出及时性代价。上图展示八个实际候选的能耗与加权延误，折线仅为这两个显示目标下的非支配关系，星号标记所选方案；下图按实体无人机展示该方案的占用区间。有限候选集不构成全局前沿证明。}{optimization}

\subsection{结果及其含义}
\textbf{所选方案用 \QTwoTrips{} 架次完成全部货箱，运输能耗 \QTwoEnergy{} kWh，完工历时 \QTwoTime{} s。}加权延误为 \QTwoLate{} 优先系数·s，全部硬截止要求满足。逐箱交付表将到达与完成交接严格区分，资源检查也覆盖电池返航后的充电阶段。

图中的实体无人机占用并不均衡：大载荷机型承担较多后续任务，小型机较早释放。这是当前候选下的任务分配结果，不应把长空闲直接归因于算法错误；但它提示还可扩大跨机型交换和离散路线搜索，以评估剩余改进空间。本文没有该问题的全局下界，故不报告未经证实的最优性差距。
\FloatBarrier
